\documentclass[11pt,a4paper]{article}

\usepackage[T1]{fontenc}
\usepackage[utf8]{inputenc}

% Stabilize PDF metadata in addition to SOURCE_DATE_EPOCH from build.py.
\ifdefined\pdfinfoomitdate\pdfinfoomitdate=1\fi
\ifdefined\pdftrailerid\pdftrailerid{}\fi
\ifdefined\pdfsuppressptexinfo\pdfsuppressptexinfo=-1\fi

\setlength{\oddsidemargin}{0mm}
\setlength{\evensidemargin}{0mm}
\setlength{\textwidth}{159mm}
\setlength{\topmargin}{-10mm}
\setlength{\textheight}{235mm}
\setlength{\parindent}{0pt}
\setlength{\parskip}{0.7em}
\raggedbottom
\sloppy

\newcommand{\code}[1]{\texttt{#1}}
\newcommand{\PipelineObserve}{Projekt\\beobachten}
\newcommand{\PipelinePlan}{Differenz\\planen}
\newcommand{\PipelineStage}{isoliert\\stagen}
\newcommand{\PipelineVerify}{Ergebnis\\prüfen}
\newcommand{\PipelineFinish}{Commit oder\\Rollback}
\newcommand{\PipelineFailure}{Fehler: unverändert zurück}
\newcommand{\ProjectBoundaryTitle}{Projektbesitz}
\newcommand{\ProjectBoundaryOne}{Anwendungsquellen}
\newcommand{\ProjectBoundaryTwo}{Domänenkonfiguration}
\newcommand{\ProjectBoundaryThree}{unbekannte Dateien}
\newcommand{\ProjectBoundaryFour}{Produkt-Repositories}
\newcommand{\ToolingBoundaryTitle}{Toolingbesitz}
\newcommand{\ToolingBoundaryOne}{Profile und Regeln}
\newcommand{\ToolingBoundaryTwo}{Adapter und Ressourcen}
\newcommand{\ToolingBoundaryThree}{Migrationen und State}
\newcommand{\ToolingBoundaryFour}{Tooling-Dokumentation}
\newcommand{\BoundaryObserve}{Evidenz}
\newcommand{\BoundaryProposal}{begrenzter Patch}
\newcommand{\BoundaryContract}{expliziter Eigentumsvertrag}

\title{Portable Toolingarchitektur\\Eine bilinguale Fallstudie zu kontrollierter Integration}
\author{Fallstudie des portablen Toolings}
\date{Reproduzierbare Ausgabe}

\begin{document}
\maketitle

\begin{abstract}
Diese Fallstudie untersucht ein Toolingpaket, das als direkt kopierbarer Ordner in
unterschiedliche Softwareprojekte integriert werden kann. Anders als ein
Produkt-Template besitzt es weder die Anwendung noch deren Lebenszyklus. Es beobachtet
vorhandene Technologien, wählt ein Profil, plant eng begrenzte Änderungen und führt diese
transaktional aus. Bewertet werden Portabilität, Eigentumsschutz, deterministische
Ausführung und die Kosten eines bewusst fehlerschließenden Designs. Die Ergebnisse zeigen,
dass die stärksten Garantien nicht aus einer großen Zahl von Generatoren entstehen, sondern
aus expliziten Grenzen, wiederholbaren Plänen und prüfbarer Rücknahme.
\end{abstract}

\tableofcontents
\newpage

\section{Ausgangslage und Forschungsfrage}

Viele interne Tooling-Repositories beginnen als vollständige Projektvorlage. Mit der Zeit
enthalten sie Annahmen über Frontend, Backend, Datenbank, Desktop-Hülle, CI und Release.
Solange jedes Zielprojekt dieselbe Form besitzt, ist das bequem. Sobald das Tooling jedoch
in ein bestehendes oder anders strukturiertes Projekt kopiert wird, werden diese Annahmen zu
Risiken: feste Pfade zeigen auf falsche Dateien, ein Generator überschreibt Produktcode und
ein Update kann nicht mehr zwischen eigenem und fremdem Zustand unterscheiden.

Der untersuchte Umbau ersetzt dieses Modell durch einen portablen Toolingordner. Seine
zentrale Forschungsfrage lautet:

\begin{quote}
Wie kann ein kopierbares Tooling verschiedene Projektformen unterstützen, ohne
Anwendungsquellen zu besitzen, und dabei Integration, Migration sowie Verifikation
reproduzierbar und rücknehmbar machen?
\end{quote}

Die Fallstudie betrachtet diese Frage nicht als reine Paketierungsaufgabe. Portabilität
bedeutet hier, dass Pfade aus einem zur Laufzeit ermittelten Projektkontext stammen, dass
Technologiekenntnis hinter Adaptern liegt und dass jede schreibende Operation vorab als
strukturierter Plan sichtbar wird. Ein kopierter Ordner ist erst dann wirklich portabel,
wenn sein Verhalten nicht von der Position des ursprünglichen Repositories abhängt.

\section{Architektur des portablen Kerns}

\subsection{Kontext statt globaler Wurzel}

Der Projektkontext trennt Projektwurzel, Toolingwurzel, Dokumentationswurzel,
Ressourcenpfade und lokalen State. Standardwurzeln werden zentral aus dem Installationsort
des Toolings abgeleitet, an der Grenze validiert und danach explizit weitergereicht.
Verbraucher erraten dadurch nicht jeweils eine eigene globale Repository-Wurzel. Dieser
kleine Unterschied ist entscheidend: Dieselbe Kopie kann unter einem Monorepo, neben einer
Webanwendung oder in einem Desktopprojekt liegen, ohne dass Quellpfade neu kompiliert werden
müssen.

Die Projektkonfiguration beschreibt Pfade und Profilwünsche, bleibt aber Projekteigentum.
Profile, Adapterverträge und kanonische Ressourcen liegen dagegen im Tooling. Der lokale
Integrationszustand befindet sich außerhalb des Toolingordners in einem begrenzten
State-Bereich. Dadurch kann der komplette Toolingordner ausgetauscht werden, ohne den
Projektzustand mit einer alten Implementierung zu vermischen.

\subsection{Profile und Adapter}

Ein Profil formuliert gewünschte Fähigkeiten, nicht eine vollständige Verzeichnisstruktur.
Beispiele sind eine reine Weboberfläche, eine Cloud-Variante oder eine Desktopkombination.
Adapter übersetzen diese Fähigkeiten in technologiespezifische Beobachtungen und
kontrollierte Aktionen. Ein Frontend-Adapter kann etwa vorhandene Paketmanifeste erkennen,
während ein Datenbank-Adapter nur bei belastbaren Markern aktiv wird. Fehlt Evidenz, bleibt
der Adapter inaktiv; bloße Vermutung reicht nicht für einen Schreibzugriff.

Diese Aufteilung hält den Kern klein. Er kennt Zustände, Findings, Pläne, Transaktionen und
Migrationen. Er kennt jedoch weder konkrete npm-Skripte noch Datenbankdateien als universelle
Wahrheit. Neue Technologien erweitern die Adapterregistrierung, ohne den Kern auf ein neues
Produktmodell festzulegen.

\section{Eigentum als Sicherheitsgrenze}

Abbildung~\ref{fig:ownership} zeigt die wichtigste Grenze des Entwurfs. Das Projekt besitzt
seine Anwendung, domänenspezifische Konfiguration und alle unbekannten Dateien. Das Tooling
besitzt ausschließlich seine mitgelieferten Regeln, Ressourcen, Migrationen und
Dokumentation. Ein Adapter darf Projektdateien beobachten und eng begrenzte Änderungen
vorschlagen; daraus folgt kein pauschales Eigentum am Zielbaum.

\begin{figure}[htbp]
  \centering
  % Portable, package-free LaTeX picture. Labels are supplied by each language source.
\begingroup
\setlength{\unitlength}{1mm}
\begin{picture}(118,54)
  \put(0,7){\framebox(50,40){%
    \parbox{43mm}{\centering\small\textbf{\ProjectBoundaryTitle}\\[2mm]
      \scriptsize\ProjectBoundaryOne\\\ProjectBoundaryTwo\\
      \ProjectBoundaryThree\\\ProjectBoundaryFour}}}
  \put(68,7){\framebox(50,40){%
    \parbox{43mm}{\centering\small\textbf{\ToolingBoundaryTitle}\\[2mm]
      \scriptsize\ToolingBoundaryOne\\\ToolingBoundaryTwo\\
      \ToolingBoundaryThree\\\ToolingBoundaryFour}}}
  \put(50,31){\vector(1,0){18}}
  \put(59,35){\makebox(0,0){\scriptsize\BoundaryObserve}}
  \put(68,20){\vector(-1,0){18}}
  \put(59,16){\makebox(0,0){\scriptsize\BoundaryProposal}}
  \put(59,51){\makebox(0,0){\small\textbf{\BoundaryContract}}}
\end{picture}
\endgroup

  \caption{Evidenz fließt zum Tooling, Änderungen nur als begrenzte Vorschläge zurück.}
  \label{fig:ownership}
\end{figure}

Der Eigentumsvertrag wirkt auf drei Ebenen:

\begin{enumerate}
  \item \textbf{Pfadgrenzen:} relative Pfade werden normalisiert, müssen innerhalb der
  Projektwurzel liegen und dürfen keine symbolischen Umleitungen ausnutzen.
  \item \textbf{Änderungsgrenzen:} strukturierte Operationen unterscheiden Erstellen,
  Ersetzen und Entfernen. Vorbedingungen verhindern, dass ein inzwischen veränderter Inhalt
  stillschweigend überschrieben wird.
  \item \textbf{Wissensgrenzen:} Adapter müssen ihre Marker und Zielpfade deklarieren.
  Unbekannte Dateien sind kein Fehler und werden nicht als aufzuräumender Rest behandelt.
\end{enumerate}

Diese Regeln sind strenger als ein klassischer Generator. Sie reduzieren den Umfang
automatischer Reparaturen, erhöhen aber die Nachvollziehbarkeit. Für ein Tool, das in fremden
Repositories laufen soll, ist diese Asymmetrie sinnvoll: Ein ausgelassener Komfortschritt
ist korrigierbar, ein verlorener Produktbestand möglicherweise nicht.

\section{Transaktionale Integration}

Die Integrationspipeline in Abbildung~\ref{fig:pipeline} trennt Beobachtung und Schreiben.
Zunächst entsteht aus Ist-Zustand, Profil und Adapterbefunden ein deterministischer Plan.
Die geplanten Operationen werden anschließend in einer isolierten Staging-Kopie angewendet.
Wenn der Plan sie erfordert, laufen feste Tooling-, Qualitäts-, Test- und
Lockfile-Abhängigkeitsprüfungen gegen diesen Kandidaten. Erst ein vollständig erfolgreicher
Kandidat darf in das Projekt übernommen werden.

\begin{figure}[htbp]
  \centering
  % Portable, package-free LaTeX picture. Labels are supplied by each language source.
\begingroup
\setlength{\unitlength}{1mm}
\begin{picture}(118,31)
  \put(0,12){\framebox(20,12){\scriptsize\shortstack{\PipelineObserve}}}
  \put(20,18){\vector(1,0){4}}
  \put(24,12){\framebox(20,12){\scriptsize\shortstack{\PipelinePlan}}}
  \put(44,18){\vector(1,0){4}}
  \put(48,12){\framebox(20,12){\scriptsize\shortstack{\PipelineStage}}}
  \put(68,18){\vector(1,0){4}}
  \put(72,12){\framebox(20,12){\scriptsize\shortstack{\PipelineVerify}}}
  \put(92,18){\vector(1,0){4}}
  \put(96,12){\framebox(22,12){\scriptsize\shortstack{\PipelineFinish}}}
  \put(83,10){\vector(0,-1){5}}
  \put(83,2){\makebox(0,0){\scriptsize\PipelineFailure}}
  \put(83,5){\line(-1,0){59}}
  \put(24,5){\vector(0,1){7}}
\end{picture}
\endgroup

  \caption{Die Integration veröffentlicht nur einen vollständig verifizierten Kandidaten.}
  \label{fig:pipeline}
\end{figure}

Der Commit ist selbst wieder kontrolliert. Vor dem Schreiben werden die beobachteten
Vorbedingungen erneut geprüft. Backups und ein Transaktionsjournal ermöglichen Rollback,
wenn eine spätere Operation fehlschlägt. Der persistierte State speichert keine pauschale
Momentaufnahme des gesamten Projekts, sondern die für Besitz und Drift notwendigen Belege.
Damit bleibt ein manuell veränderter Projektbestand sichtbar, statt beim nächsten Lauf als
veraltete Vorlage ersetzt zu werden.

Zwei Betriebsarten ergeben sich aus demselben Planungsmodell. Eine reine Prüfung erzeugt
Findings und einen Plan, verändert aber kein Byte. Eine vollständige Reparatur darf nur die
im Plan autorisierten Operationen ausführen. Nach erfolgreichem Abschluss muss eine erneute
Prüfung keinen Drift mehr melden; ein zweiter Reparaturlauf muss ein No-op sein. Diese
Idempotenz ist ein stärkeres Akzeptanzkriterium als der einmalige Exitcode eines Generators.

\section{Ausführung fremder Werkzeuge}

Manche Verifikation benötigt reale Toolchains. Ein Staging-Kandidat kann vertrauenswürdige
Tooling-Qualitäts- und Testprüfungen oder eine Lockfile-Validierung über ein offline
ausgeführtes \code{npm ci --ignore-scripts} benötigen. Getrennt davon kann der Benutzer eine
feste Live-Capability eines vorhandenen Produkts ausdrücklich auslösen. Das Tooling kann
diesen Code nicht semantisch ersetzen, muss seine Ausführung aber begrenzen; frei
konfigurierte Shellfragmente gehören nicht zum Vertrag.

Beide Aktionspfade laufen ohne Shell, mit beschränkter Ausgabe, Timeout und temporären Home-,
Cache- und Konfigurationsverzeichnissen. Die Staging-Validierung erzwingt für unterstützte
Werkzeuge zusätzlich den Offline-Modus und deaktiviert Lifecycle-Skripte; eine explizite
Live-Aktion darf dagegen bewusst Registries kontaktieren und Projektartefakte erzeugen.
Prozessgruppen beziehungsweise Job-Objekte beenden Nachfahren nach Abschluss. Diese
Maßnahmen sind keine Betriebssystem-Sandbox: Ausgeführter Code besitzt weiterhin die Rechte
des lokalen Benutzers. Sie machen die verbleibende Vertrauensannahme explizit; nur die
Staging-Validierung gehört zur Rollback-Grenze der Integration.

\section{Evaluation}

\subsection{Versuchsaufbau}

Die Architektur wird mit fünf Projektprofilen und unabhängigen temporären Zielprojekten
bewertet. Für jedes Profil gilt dieselbe Sequenz: Tooling kopieren, Projektzustand erkennen,
prüfen, vollständig integrieren, erneut prüfen, passende Tests ausführen und die Integration
ein zweites Mal starten. Ein ergänzender Versuch integriert die fest referenzierte historische
Version 0.1.0 und ersetzt danach beide verwalteten Verzeichnisse durch das aktuelle Release,
um eine registrierte Migration und die Trennung von Tooling- und Projektstate zu prüfen.

Die Messgrößen sind bewusst binär und byteorientiert:

\begin{itemize}
  \item Hashes aller unbekannten und produktbesessenen Dateien bleiben unverändert.
  \item Eine reine Prüfung verändert weder Dateien noch Verzeichnisstruktur.
  \item Nach erfolgreicher Integration liefert die zweite Integration keine Operation.
  \item Ein absichtlich fehlschlagender Verifikationsschritt stellt den Ausgangszustand her.
  \item Dieselbe portable Nutzlast besitzt nach direkter Kopie dieselbe Manifestidentität.
\end{itemize}

\subsection{Ergebnis und Interpretation}

\begin{center}
\begin{tabular}{|p{43mm}|p{48mm}|p{55mm}|}
\hline
\textbf{Kriterium} & \textbf{Beobachtbarer Nachweis} & \textbf{Interpretation} \\
\hline
Relokation & Kontextpfade zeigen nur auf das temporäre Ziel & Keine Bindung an das
Quellrepository \\
\hline
Eigentumsschutz & Produkt- und unbekannte Hashes bleiben gleich & Adapterwissen begrenzt
statt erweitert Besitz \\
\hline
Idempotenz & Zweiter vollständiger Lauf plant null Operationen & State und gewünschter
Zustand stimmen überein \\
\hline
Rollback & Fehlerkandidat hinterlässt keine Teiländerung & Verifikation liegt vor
Veröffentlichung \\
\hline
Toolingtausch & Migration läuft gegen erhaltenen Projektstate & Implementierung und State
sind getrennt \\
\hline
\end{tabular}
\end{center}

Die Evaluation belegt nicht, dass jeder denkbare Adapter korrekt ist. Sie belegt die
Architektureigenschaften, die für alle Adapter gelten müssen. Der Transaktionskern kann
Ausführungs- und Verifikationsfehler zurückrollen, aber nicht beweisen, dass eine formal gültige
Änderung an einem erlaubten Domänenschlüssel fachlich erwünscht ist. Umgekehrt kann ein
inhaltlich guter Adapter die Portabilität verletzen, wenn er feste Wurzelpfade verwendet.
Darum kombinieren die Tests Profilmatrizen mit Vertragsprüfungen für Pfade, Marker,
Eigentumsbereiche und Aktionsdefinitionen.

\section{Trade-offs und Grenzen}

\textbf{Mehr Vorarbeit, kleinere Fehlerfläche.} Ein Adapter benötigt deklarierte Marker,
Operationen und Tests. Das ist aufwendiger als das Kopieren eines fertigen Verzeichnisbaums,
verhindert aber implizite Besitzannahmen.

\textbf{Fail-closed kann Verfügbarkeit kosten.} Fehlt ein Lockfile, ein Compiler oder eine
sichere Pfadauflösung, verweigert das Tooling die automatische Änderung. Teams müssen solche
Befunde bewusst beheben. Der Vorteil ist, dass Unsicherheit nicht als Erlaubnis interpretiert
wird.

\textbf{Reproduzierbarkeit ist umgebungsgebunden.} Fixierte Zeit, Locale, Offline-Schalter
und bereinigte Caches beseitigen viele Quellen von Drift. Unterschiedliche Compiler- oder
Toolchainversionen können trotzdem andere Artefakte erzeugen. Vollständige
Byte-Reproduzierbarkeit benötigt daher zusätzlich eine versionierte Buildumgebung.

\textbf{Prozessbegrenzung ist keine Sandbox.} Timeouts und Prozessbäume verhindern typische
Lecks, isolieren aber nicht die Dateisystemrechte eines bösartigen Tests. Für nicht
vertrauenswürdige Projekte ist eine externe VM- oder Containergrenze erforderlich.

\textbf{Das Manifest ist eine Grenze der Selbstkonsistenz.} Eine exakte Nutzlastliste erkennt
gemischte oder veränderte Toolingkopien nur, solange das mitgelieferte Manifest als
vertrauenswürdig gilt. Koordinierte Änderungen an Nutzlast und Manifest werden nicht
authentifiziert. Nach jeder legitimen Änderung muss es neu erzeugt werden; ein veraltetes
Manifest blockiert portable Integration und Verifikation, während ein Release-Gate diese
Prüfung ausdrücklich aufrufen muss.

\section{Schlussfolgerung}

Der Wechsel vom Master-Template zum portablen Tooling ist vor allem ein Wechsel des
Eigentumsmodells. Profile beschreiben Fähigkeiten, Adapter liefern begrenztes
Technologiewissen und der Kern veröffentlicht nur verifizierte Transaktionen. Portabilität
entsteht durch expliziten Kontext; Sicherheit entsteht durch enge Pfade und fehlerschließende
Entscheidungen; Wartbarkeit entsteht durch die Trennung von Projektstate und austauschbarer
Toolingimplementierung.

Die Fallstudie zeigt zugleich die Grenze dieser Aussage. Das System kann fremden Code nicht
magisch vertrauenswürdig machen und Toolchainversionen nicht ohne Buildumgebung vereinheitlichen.
Es kann aber jede dieser Annahmen sichtbar machen, den Änderungsumfang beweisen und bei Fehlern
zum bekannten Ausgangszustand zurückkehren. Für ein Tool, das in vielen unterschiedlich
gewachsenen Projekten arbeiten soll, ist genau diese kontrollierte Bescheidenheit die
entscheidende Architekturleistung.

\end{document}
